% !TeX root = ../main.tex

\documentclass[../main.tex]{subfiles}
\begin{document}

\chapter{CƠ SỞ LÝ THUYẾT VÀ THUYẾT MINH QUY TRÌNH CÔNG NGHỆ}

\section{Cơ sở lựa chọn công nghệ nhiệt phân cho dự án Hội An}

Như đã phân tích ở Chương 1, dự án xử lý chất thải rắn sinh hoạt tại thành phố Hội An ban đầu được phê duyệt sử dụng công nghệ khí hóa kết hợp tuabin khí. Tuy nhiên, qua quá trình triển khai và khảo sát thực tế, phương án này bộc lộ nhiều hạn chế không phù hợp với đặc thù rác thải địa phương. Rác thải sinh hoạt tại Hội An có độ ẩm rất cao (dao động từ 50 -- 65\%) và thành phần phức tạp, biến động theo mùa vụ du lịch. Trong khi đó, công nghệ khí hóa đòi hỏi nguyên liệu đầu vào phải được sấy khô đến độ ẩm dưới 20 -- 25\% và có độ đồng nhất cao, dẫn đến việc tiêu tốn một lượng năng lượng khổng lồ cho công đoạn tiền xử lý, làm giảm hiệu quả kinh tế chung của toàn bộ hệ thống.

Bên cạnh đó, khí tổng hợp (syngas) sinh ra từ quá trình khí hóa rác thải thường chứa nhiều tạp chất như hắc ín (tar), hạt bụi mịn và các hợp chất chứa lưu huỳnh, clo. Việc làm sạch dòng khí này để đáp ứng tiêu chuẩn khắt khe của tuabin khí đòi hỏi hệ thống xử lý cực kỳ phức tạp và đắt tiền. Bài học từ các dự án khí hóa rác thải quy mô thử nghiệm tại Việt Nam cho thấy sự không ổn định trong chất lượng khí đầu ra thường xuyên gây tắc nghẽn và hư hỏng thiết bị, dẫn đến hiệu suất vận hành thấp và chi phí bảo trì tăng vọt.

Đứng trước những thách thức đó, việc chuyển đổi sang công nghệ nhiệt phân (pyrolysis) được xem là một bước đi đột phá và mang tính thực tiễn cao. Công nghệ nhiệt phân thể hiện sự linh hoạt vượt trội khi có thể xử lý hiệu quả nguồn nguyên liệu rác nhựa đã qua phân loại (có độ ẩm thấp) mà không đòi hỏi sự đồng nhất tuyệt đối về thành phần. Điểm khác biệt cốt lõi là quá trình nhiệt phân tạo ra sản phẩm chính là dầu lỏng (bio-oil). Loại nhiên liệu này có thể được lưu trữ, vận chuyển dễ dàng và sử dụng linh hoạt cho các mục đích khác nhau như chạy máy phát điện diesel hoặc gia nhiệt cho chính lò phản ứng, tạo nên một chu trình năng lượng khép kín. Việc sử dụng lò nhiệt phân quay cấu trúc mô-đun cũng cho phép nhà máy dễ dàng mở rộng công suất trong tương lai bằng cách lắp đặt thêm các đơn vị song song mà không làm gián đoạn dây chuyền hiện hữu.

\section{Cơ sở lý thuyết về công nghệ nhiệt phân}

\subsection{Nguyên lý quá trình nhiệt phân}

Nhiệt phân (Pyrolysis) là một quá trình biến đổi hóa học nhiệt phân hủy các vật liệu hữu cơ trong môi trường hoàn toàn thiếu oxy hoặc có nồng độ oxy rất thấp. Khác với quá trình đốt cháy trực tiếp (incineration) -- nơi vật liệu phản ứng mạnh với oxy để tạo ra nhiệt lượng và các khí nhà kính (\ce{CO2}, \ce{H2O}), nhiệt phân là quá trình cracking nhiệt (thermal cracking). Dưới tác dụng của nhiệt lượng cao, các liên kết hóa học phức tạp và chuỗi polymer dài trong cấu trúc của rác thải (đặc biệt là rác nhựa) bị bẻ gãy một cách ngẫu nhiên hoặc ở cuối chuỗi, tạo thành các phân tử có cấu trúc ngắn hơn và đơn giản hơn.

Đối với rác nhựa, vốn là thành phần chủ đạo cấp cho lò nhiệt phân tại nhà máy Hội An, quá trình này diễn biến qua ba giai đoạn nhiệt động học chính. Giai đoạn đầu tiên là sấy và làm mềm diễn ra ở dải nhiệt độ dưới 200°C, tại đây lượng ẩm dư thừa và các khí hòa tan bay hơi hoàn toàn, khối nhựa bắt đầu chuyển sang trạng thái mềm dẻo và tan chảy. Tiếp đó là giai đoạn nhiệt phân sơ cấp (200 -- 350°C), các liên kết hóa học cốt lõi như C--C và C--H bắt đầu bị đứt gãy, sản sinh ra phần lớn các hydrocarbon có nhiệt độ sôi cao (dầu nặng) và các oligomer trung gian. Cuối cùng, tại giai đoạn nhiệt phân thứ cấp (350 -- 500°C), các sản phẩm sơ cấp tiếp tục chịu tác động của nhiệt độ cao, bị phân cắt thành các hydrocarbon nhẹ hơn, đồng thời giải phóng một lượng lớn khí không ngưng (syngas) và để lại phần cặn rắn chứa carbon (biochar).

\subsection{Sản phẩm của quá trình nhiệt phân}

Sự phân bố các sản phẩm nhiệt phân phụ thuộc mạnh mẽ vào nhiệt độ, tốc độ gia nhiệt và thời gian lưu của vật liệu trong lò. Nhìn chung, quá trình này tạo ra đồng thời ba nhóm sản phẩm chính, mỗi nhóm đều có giá trị ứng dụng riêng biệt trong mô hình kinh tế tuần hoàn.

Sản phẩm giá trị nhất là \textbf{dầu nhiệt phân (Bio-oil)}. Đây là một hỗn hợp chất lỏng phức tạp gồm các hydrocarbon mạch thẳng và mạch vòng, thu được từ việc ngưng tụ hơi thoát ra khỏi lò. Tùy thuộc vào chất lượng nguyên liệu, tỷ lệ thu hồi dầu có thể dao động từ 12\% đến 50\% khối lượng. Dầu nhiệt phân có nhiệt trị rất cao, trung bình đạt khoảng 43,26 MJ/kg, tương đương với nhiên liệu diesel thương phẩm, biến nó thành nguồn năng lượng lý tưởng cho các tổ máy phát điện.

Sản phẩm thứ hai là \textbf{khí tổng hợp không ngưng (Syngas)}. Dòng khí này bao gồm chủ yếu là methane (\ce{CH4}), ethylene (\ce{C2H4}), hydrogen (\ce{H2}), carbon monoxide (\ce{CO}) và một lượng nhỏ carbon dioxide (\ce{CO2}). Mặc dù có nhiệt trị thấp hơn dầu lỏng, syngas đóng vai trò then chốt trong việc tối ưu hóa năng lượng của nhà máy. Khí này được thu hồi và dẫn ngược trở lại buồng đốt ngoài của lò nhiệt phân, cung cấp lượng nhiệt đáng kể để duy trì quá trình phản ứng, từ đó giảm thiểu tối đa sự phụ thuộc vào nguồn nhiên liệu đốt bên ngoài.

Sản phẩm cuối cùng là \textbf{than sinh học (Biochar)}, phần chất rắn còn sót lại ở đáy lò chiếm khoảng 20 -- 25\% khối lượng nguyên liệu đầu vào. Biochar chứa cấu trúc carbon cố định và một số khoáng chất trơ, có đặc tính xốp và khả năng giữ nước tốt. Sản phẩm này có tiềm năng ứng dụng rất lớn trong lĩnh vực nông nghiệp, đóng vai trò như một chất cải tạo đất, giúp cải thiện độ màu mỡ hoặc có thể được nâng cấp thành vật liệu hấp phụ (than hoạt tính) trong xử lý môi trường.

\section{Sơ đồ quy trình công nghệ tổng thể}

Nhà máy xử lý chất thải rắn sinh hoạt thành phố Hội An được thiết kế và vận hành theo mô hình "điện rác" hiện đại và khép kín. Quy trình công nghệ là sự kết hợp liên hoàn giữa nhiều khối chức năng, đảm bảo dòng vật chất và dòng năng lượng được lưu chuyển một cách tối ưu, giảm thiểu phát thải thứ cấp ra môi trường. Cấu trúc hệ thống bao gồm bốn khối công nghệ chính: tiền xử lý và sản xuất viên nén RDF; nhiệt phân và carbon hóa; tinh chế pha trộn nhiên liệu; và khối phát điện.

Sự liên kết này thể hiện tính tuần hoàn nội tại ở mức độ cao. Phế phẩm của công đoạn này ngay lập tức trở thành nguyên liệu hoặc nguồn năng lượng cho công đoạn khác. Chẳng hạn, nhiệt dư từ hệ thống xử lý khí thải được tận dụng để sấy rác, khí syngas duy trì nhiệt cho lò phản ứng, và nước ngưng được tuần hoàn làm mát.

\begin{figure}[H]
\centering
% Vị trí chèn ảnh sơ đồ công nghệ tổng quát
\fbox{\parbox{0.9\textwidth}{\centering \vspace{3cm} \textit{Sơ đồ dây chuyền công nghệ xử lý rác và phát điện tại Hội An (Sinh viên bổ sung hình vẽ tại đây)} \vspace{3cm}}}
\caption{Sơ đồ luồng vật chất và năng lượng tổng thể của nhà máy}
\label{fig:sodoQTCN_ChiTiet}
\end{figure}

\section{Thuyết minh chi tiết các công đoạn công nghệ}

\subsection{Hệ thống tạo viên nén nhiên liệu RDF (Khối 1)}

Hệ thống sản xuất viên nén nhiên liệu có nguồn gốc từ rác thải (RDF -- Refuse Derived Fuel) đóng vai trò tiền xử lý thiết yếu, chuyển hóa dòng rác thải sinh hoạt đầu vào hỗn độn thành nguồn nguyên liệu đồng nhất phục vụ cho lò nhiệt phân. Quy trình này bao gồm chuỗi các hoạt động cơ lý hóa được kiểm soát chặt chẽ.

Ban đầu, rác thải tiếp nhận tại nhà máy được phun chế phẩm vi sinh EM (Effective Microorganisms) nhằm ức chế hoạt động của vi khuẩn gây bệnh và khử triệt để mùi hôi khó chịu. Sau đó, rác được đưa qua máy xé bao và hệ thống phân loại kết hợp cơ khí -- thủ công. Tại đây, kim loại được tách ra nhờ nam châm từ tính, các thành phần trơ (như thủy tinh, xà bần) và chất thải nguy hại (như pin, hóa chất) được phân lập để xử lý theo quy định riêng. Phần rác cháy được (bao gồm nhựa, giấy, nilon, hữu cơ) tiếp tục được đưa vào máy cắt băm nhỏ đến kích thước đồng đều từ 2,5 -- 5,0 cm.

Một bước đột phá trong hệ thống là quá trình bức xạ "sốc nhiệt" ở 200°C trong thời gian cực ngắn. Bước này sử dụng năng lượng điện để tiêu diệt hoàn toàn vi sinh vật, đình chỉ mọi quá trình phân hủy sinh học tự nhiên, đồng thời hạ một phần độ ẩm của rác. Sau bức xạ, rác được nén ép với áp suất cao để vắt kiệt nước, đưa độ ẩm xuống còn khoảng 25 -- 30\%. Lượng rác này tiếp tục được sấy khô bằng nhiệt dư thu hồi từ hệ thống khói thải của lò nhiệt phân cho đến khi độ ẩm đạt ngưỡng tối ưu 10 -- 16\%. Cuối cùng, hỗn hợp nguyên liệu được đưa qua máy ép trục ngang để tạo thành các viên nén RDF có khối lượng riêng cao, nhiệt trị ổn định (đạt khoảng 4.600 kcal/kg) và được lưu trữ trong silo kín, sẵn sàng cấp cho lò phản ứng.

\subsection{Hệ thống nhiệt phân nhiên liệu RDF (Khối 2)}

Nguyên liệu chính đưa vào lò nhiệt phân chủ yếu là thành phần rác nhựa có trong viên nén RDF. Thiết bị trung tâm của hệ thống là lò nhiệt phân quay (rotary kiln) dạng xi lanh nằm ngang, vỏ thép chịu nhiệt, quay với tốc độ chậm (0,5 -- 2 vòng/phút) để đảo trộn đều vật liệu bên trong.

Quá trình nhiệt phân được thiết kế theo chế độ vận hành từng mẻ (batch). Sau khi nạp đủ lượng rác nhựa (khoảng 6 -- 10 tấn/mẻ), hệ thống được niêm phong hoàn toàn để đảm bảo điều kiện yếm khí. Đầu đốt ngoài bắt đầu cấp nhiệt, nâng nhiệt độ buồng lò từ từ vượt qua ngưỡng 190°C, lúc này các liên kết polymer bắt đầu bẻ gãy. Giai đoạn phản ứng mạnh nhất diễn ra khi nhiệt độ duy trì ổn định trong khoảng 300 -- 345°C. Hơi hydrocarbon liên tục thoát ra khỏi buồng phản ứng và được hút về hệ thống ngưng tụ đa cấp. Khí syngas sinh ra được lọc sạch các hợp chất lưu huỳnh sơ bộ rồi dẫn ngược lại đầu đốt để duy trì nhiệt, nhờ đó tiết kiệm đến 80\% lượng nhiên liệu mồi bên ngoài. Sau chu kỳ khoảng 20 -- 24 giờ, mẻ phản ứng kết thúc, lò được làm nguội và tiến hành xả phần than biochar ra khỏi đáy lò thông qua hệ thống băng tải kín.

\begin{table}[H]
\centering
\caption{Thông số kỹ thuật vận hành lò nhiệt phân}
\label{tab:ts_nhietphan}
\renewcommand{\arraystretch}{1.3}
\begin{tabular}{|l|c|}
\hline
\textbf{Thông số thiết kế} & \textbf{Giá trị quy định} \\
\hline
Chế độ vận hành & Theo mẻ (batch) \\
\hline
Khối lượng cấp liệu mỗi mẻ & 6 -- 10 tấn (tùy thành phần nhựa) \\
\hline
Dải nhiệt độ phản ứng tối ưu & 300 -- 350°C (tối đa có thể lên 600°C) \\
\hline
Thời gian hoàn thành một mẻ & 20 -- 24 giờ \\
\hline
Sản lượng dầu thu hồi trung bình & 166 -- 200 lít/tấn nguyên liệu \\
\hline
\end{tabular}
\end{table}

\subsection{Công nghệ carbon hóa xử lý rác hữu cơ}

Đối với phần rác thải hữu cơ (thức ăn thừa, rau củ quả) chiếm tỷ trọng lớn trong rác sinh hoạt Hội An, nhà máy ứng dụng công nghệ carbon hóa liên tục. Khác với nhiệt phân nhựa, quá trình carbon hóa đòi hỏi dải nhiệt độ vận hành cao hơn, thường từ 450 -- 590°C trong môi trường thiếu oxy. Quá trình này biến đổi triệt để các cấu trúc sinh khối phức tạp thành than sinh học và khí gas.

Rác hữu cơ sau khi phân loại có độ ẩm rất cao nên bắt buộc phải trải qua thiết bị trống sấy trước khi tiến vào lò phản ứng. Nguồn nhiệt cho trống sấy được lấy trực tiếp từ phần nhiệt thải thừa của chính lò carbon hóa, giúp hệ thống đạt hiệu suất sử dụng năng lượng tối đa. Than biochar thu được từ rác hữu cơ có chất lượng xốp, ổn định, và sau khi kiểm định đáp ứng các tiêu chuẩn an toàn về kim loại nặng, sẽ được định hướng sử dụng làm phân bón hữu cơ hoặc chất cải tạo đất phục vụ nền nông nghiệp sinh thái tại địa phương.

\subsection{Hệ thống tinh chế và chưng cất dầu nhiệt phân (Khối 3)}

Dầu nhiệt phân thô ngưng tụ trực tiếp từ lò thường chứa một lượng đáng kể nước nhũ hóa, cặn cơ học, axit hữu cơ tự do và các phân đoạn hydrocarbon mạch rất dài. Với đặc tính độ nhớt cao, chỉ số cetane thấp và hàm lượng lưu huỳnh chưa đạt chuẩn, loại dầu thô này không thể bơm trực tiếp vào động cơ diesel hiện đại. Do đó, hệ thống tinh chế đa bậc (Three-stage Refining) được thiết lập nhằm nâng cao phẩm cấp của dầu.

\begin{itemize}
    \item \textbf{Bước 1 -- Tiền xử lý xúc tác lỏng:} Dầu thô được bơm vào bể phản ứng kết hợp khuấy trộn với chất xúc tác axit dạng lỏng. Tác nhân này đóng vai trò phá vỡ bước đầu các phân tử hydrocarbon siêu nặng, làm giảm độ nhớt và phân tách sơ bộ lượng ẩm cùng tạp chất lơ lửng.
    \item \textbf{Bước 2 -- Cracking xúc tác rắn:} Dòng dầu tiếp tục được gia nhiệt hóa hơi và đẩy qua tháp chứa các tầng vật liệu xúc tác rắn (thường là zeolite hoặc alumina). Tại đây diễn ra hàng loạt các phản ứng cracking thứ cấp, bẻ gãy mạnh mẽ các chuỗi carbon dài ($C_{15} - C_{30}$) thành các phân đoạn tương đương với dải điểm sôi của dầu diesel và xăng ($C_8 - C_{15}$).
    \item \textbf{Bước 3 -- Keo tụ và làm trong:} Sản phẩm lỏng sau khi ngưng tụ lại từ tháp cracking được xử lý bằng hóa chất keo tụ (flocculant). Các cặn bẩn phân cực và nước vi nhũ tương sẽ liên kết lại thành khối lớn và lắng xuống đáy, trả lại lớp dầu trong suốt, nhẹ và có phẩm chất cao hơn nhiều so với dầu thô ban đầu.
\end{itemize}

Mặc dù đã trải qua chưng cất, dầu thành phẩm vẫn tồn tại một vài chỉ tiêu chưa đạt chuẩn hoàn hảo (như chỉ số cetane và điểm chớp cháy). Vì vậy, giải pháp tối ưu được áp dụng tại nhà máy là \textit{pha trộn nhiên liệu} (blending). Dầu nhiệt phân tinh chế được hòa trộn tự động với nhiên liệu diesel thương mại (ULSD) theo tỷ lệ 30\% dầu sinh học và 70\% diesel. Ở tỷ lệ vàng này, hỗn hợp nhiên liệu mới hoàn toàn đáp ứng các thông số vận hành an toàn cho động cơ mà không đòi hỏi bất kỳ sự cải hoán phức tạp nào về cấu trúc máy phát.

\subsection{Hệ thống phát điện động cơ diesel (Khối 4)}

Sau khi đánh giá và loại bỏ các phương án dùng tuabin khí (do không đáp ứng yêu cầu chất lượng dòng khí syngas) và tuabin hơi nước (do chi phí đầu tư quá lớn và hiệu suất điện thấp, chỉ đạt khoảng 14\%), phương án tổ máy phát điện chạy bằng động cơ diesel lai hợp nhiên liệu đã được lựa chọn.

Tổ hợp máy phát điện bao gồm một động cơ đốt trong 4 kỳ đồng bộ với máy phát điện xoay chiều có công suất định mức 500 kW. Động cơ hoạt động dựa trên nguyên lý tự cháy do nén nổ: hỗn hợp nhiên liệu (bio-oil pha diesel) được bơm cao áp phun tơi vào buồng đốt dưới áp suất cực lớn vào thời điểm cuối kỳ nén, khi nhiệt độ không khí đã lên tới 550 -- 650°C. Phản ứng cháy giãn nở đẩy piston sinh công cơ học làm quay trục khuỷu, từ đó lai dắt rotor máy phát tạo ra dòng điện ba pha tần số 50 Hz.

Hiệu suất nhiệt của động cơ diesel rất cao, dao động từ 30 -- 40\%, đảm bảo chuyển hóa tối đa hóa năng của hỗn hợp dầu thành điện năng. Lượng điện này được cấp ngược lại cho các tủ điện phân phối để phục vụ toàn bộ hệ thống máy móc của nhà máy (tự dùng khoảng 30\%), phần điện năng dư thừa còn lại sẽ được đồng bộ pha và hòa thẳng vào lưới điện quốc gia.

\subsection{Hệ thống xử lý khói thải}

Hệ thống kiểm soát khí thải là một trong những hạng mục được đầu tư mạnh mẽ nhất nhằm đảm bảo nhà máy không gây ra bất cứ tác động xấu nào đến khu vực dân cư xung quanh. Do đặc thù phát sinh, khói thải được phân chia thành hai nhánh xử lý sơ bộ trước khi gộp chung.
\begin{itemize}
    \item \textbf{Nhánh 1:} Tiếp nhận khói từ lò sấy RDF và lò nhiệt phân. Dòng khí được làm mát bằng hệ thống dập ướt, sau đó đi qua cyclone để loại bỏ tro bụi có kích thước hạt lớn.
    \item \textbf{Nhánh 2:} Chịu trách nhiệm hút khói từ hệ thống carbon hóa và khí xả nhiệt độ cao từ động cơ diesel. Nguồn khói này đi qua thiết bị trao đổi nhiệt để thu hồi năng lượng nhiệt, sau đó cũng qua tháp cyclone làm mát và lọc bụi.
\end{itemize}

Hai nhánh khí thải sau tiền xử lý được hợp lưu và đẩy qua \textbf{Hệ thống lọc bụi tĩnh điện (ESP -- Electrostatic Precipitator)}. Tại ESP, dưới tác dụng của điện trường hàng chục nghìn Volt, các hạt bụi mịn li ti bị ion hóa, mất trung hòa điện và bị hút dính chặt vào các bản cực âm. Búa gõ định kỳ sẽ gõ rụng lớp bụi này xuống phễu thu tro bay. Hiệu suất lọc bụi của ESP lên tới hơn 99,5\%.

Dòng khí sạch bụi tiếp tục đi vào Tháp hấp thụ ướt chứa các tác nhân oxy hóa mạnh (Hypochlorite) để khử các khí độc như \ce{CO}, \ce{SO2}, \ce{NO2}. Tầng xử lý cuối cùng là tháp hấp phụ sợi thủy tinh chuyên dụng để khử lưu huỳnh. Khí xả cuối cùng đi ra ống khói đường kính 1.000 mm đạt tuyệt đối các quy chuẩn khắt khe nhất của Việt Nam về khí thải công nghiệp (QCVN 19:2009/BTNMT) và khí thải lò đốt chất thải rắn (QCVN 61-MT:2016/BTNMT).

\subsection{Hệ thống quản lý chất thải rắn và tro xỉ}

Phế thải rắn sinh ra trong quá trình vận hành nhà máy bao gồm ba thành phần chính: phần rác trơ không cháy được (từ khâu phân loại đầu vào), tro đáy/than biochar (từ lò nhiệt phân và carbon hóa), và tro bay (từ hệ thống lọc bụi ESP).

Phần rác không cháy như gạch đá, xà bần, thủy tinh vỡ (chiếm khoảng 13\% tổng lượng rác) sẽ được thu gom và đưa đi chôn lấp hợp vệ sinh theo quy định. Phần kim loại được thu hồi và bán phế liệu.

Than biochar và tro đáy nếu đạt tiêu chuẩn về hàm lượng kim loại nặng sẽ được xuất bán làm phân bón. Trong trường hợp không đạt, chúng sẽ được đóng bao an toàn. Đặc biệt đối với tro bay từ hệ thống lọc bụi tĩnh điện -- thành phần được xếp vào danh mục chất thải có nguy cơ chứa tạp chất nguy hại, nhà máy áp dụng quy trình cố định hóa (solidification). Tro bay được trộn đều với xi măng và các chất kết dính thành khối rắn vững chắc, ngăn chặn hoàn toàn khả năng phát tán bụi hoặc hòa tan các ion kim loại vào môi trường nước, trước khi được chuyển giao cho các đơn vị xử lý chất thải nguy hại chuyên nghiệp.

\end{document}
